\section{Mezcla de Expertos: teoría e implementación}
\label{app:moe}

Este anexo documenta el contrato de Mezcla de Expertos (MoE) implementado en
\texttt{src/model/moe/} y declarado en
\texttt{src/schema/\_model/\_moe.yaml}. El esquema es la fuente de verdad
ejecutable: cada clave discutida aquí existe en el esquema, se materializa en el
bloque anidado \texttt{moe} de \texttt{FrankensteinModelConfig} y está cubierta
por la suite de pruebas. Las dos primeras secciones son la rampa de entrada:
explican qué hace una capa MoE---una ecuación, una analogía, un ejemplo
numérico---y después los dos modos de falla (el colapso de ruteo y el rezagado
del experto caliente) que fuerzan casi todo el diseño. El resto es la
referencia: puntuación del router, estrategia de ruteo, balanceo de carga,
expertos, disposición por capas, capacidad y monitoreo, y cierra con las tablas
del contrato de esquema y de los presets. Quien termine la
Sección~\ref{app:moe-failures-es} debería poder abrir cualquier sección
posterior y saber cuál de las dos fallas, o cuál de las siete preguntas de
diseño, está respondiendo.

\subsection{Visión general: qué es una capa de Mezcla de Expertos}
\label{app:moe-overview-es}

La mitad feed-forward de un bloque transformer es \emph{densa}: una única red,
aplicada igual a cada token, sin importar qué sea el token. Eso es una clínica
con un solo médico que debe atender todas las consultas. Una capa MoE convierte
la clínica en una nómina de $E$ especialistas con una recepcionista: el
\emph{router} lee cada token y lo envía a los $k$ especialistas más pertinentes,
los \emph{expertos} $\{E_e\}_{e=1}^{E}$. En términos formales,
\begin{equation}
  y_t = g \sum_{j=1}^{k} \hat w_{t,j}\, E_{i_{t,j}}(x_t) + y^{\text{shared}}_t ,
  \label{eq:moe-layer-es}
\end{equation}
leída de adentro hacia afuera: el experto $i_{t,j}$ procesa el token y produce
un vector, el peso de mezcla $\hat w_{t,j}$ lo escala, las $k$ contribuciones se
suman, el factor de escala ruteado $g$ restaura la magnitud de la rama
(Sección~\ref{app:moe-router-es}) e $y^{\text{shared}}_t$ es la rama opcional
de expertos compartidos siempre activa (Sección~\ref{app:moe-experts-es}).
Solo se ejecutan los expertos seleccionados, y ese es todo el trato: los
parámetros totales crecen con $E$ mientras el cómputo por token queda
proporcional a $k$ \citep{shazeer2017moe,fedus2021switch}.

El trato en números: con $E = 8$ expertos del ancho denso y $k = 2$, la capa
almacena $8\times$ los parámetros feed-forward pero cada token lo procesan solo
$2$ expertos---el $25\%$ del pool---más un router de $d \times E$ pesos,
despreciable frente a cualquier experto. La memoria crece con $E$; los FLOPs
crecen con $k$. Esos son dos presupuestos independientes, y la mayor parte del
esquema es la maquinaria para gastarlos bien.

Siete preguntas abarcan el espacio de diseño, y la
Tabla~\ref{tab:moe-questions-es} mapea cada una a sus perillas y a su sección:
(i) cómo los logits se convierten en pesos, (ii) quién selecciona a los
expertos, (iii) cómo se mantiene la carga balanceada, (iv) cómo son los
expertos, (v) qué capas son MoE, (vi) qué pasa cuando un experto se desborda y
(vii) cómo se observa todo esto durante el entrenamiento.

\begin{table}[H]
\centering
\caption{Las siete preguntas de una capa MoE, las perillas que las responden y
la sección que desarrolla cada respuesta. Cada clave \texttt{model.moe.*} del
contrato de esquema (Tabla~\ref{tab:moe-schema-es}) pertenece a exactamente una
fila.}
\label{tab:moe-questions-es}
\small
\begin{tabular}{@{}L{2.2cm}L{4.8cm}L{4.6cm}L{1.6cm}@{}}
\toprule
\textbf{Eje} & \textbf{Pregunta} & \textbf{Perillas} & \textbf{Sección} \\
\midrule
Puntuación & ¿cuánto pesa cada experto seleccionado? & \texttt{router.*} & \ref{app:moe-router-es} \\
Ruteo & ¿quién elige a quién? & \texttt{routing.type}, \texttt{routing.hash\_*} & \ref{app:moe-routing-es} \\
Balanceo & ¿cómo mantenemos empleado a cada experto? & \texttt{balancing.*} & \ref{app:moe-balancing-es} \\
Expertos & ¿cuántos, qué ancho, compartidos o latentes? & \texttt{num\_experts}, \texttt{expert\_hidden\_size}, \texttt{shared.*}, \texttt{latent.*}, \texttt{experts.*} & \ref{app:moe-experts-es} \\
Disposición & ¿qué capas son MoE? & \texttt{layout.*} & \ref{app:moe-layout-es} \\
Capacidad & ¿qué pasa cuando un experto se desborda? & \texttt{routing.capacity\_factor} & \ref{app:moe-capacity-es} \\
Monitoreo & ¿cómo se observa todo? & escalares registrados & \ref{app:moe-monitoring-es} \\
\bottomrule
\end{tabular}
\end{table}

\subsection{Los dos modos de falla: colapso y rezagado}
\label{app:moe-failures-es}

Cuatro de las siete preguntas existen por dos modos de falla. Ambos se siguen
de una sola decisión de diseño---el router se aprende y es libre de
concentrarse---y cada mecanismo del resto del anexo (pérdida auxiliar, sesgo de
experto, elección por experto, capas hash, factor de capacidad) responde a uno
de ellos.

\paragraph{Colapso.}
Las decisiones de ruteo no están coordinadas: cada token elige de forma
independiente sus $k$ favoritos y nada impide que todos los tokens elijan al
mismo experto. El bucle se refuerza solo. Un experto que arranca levemente
adelante recibe más tokens; más tokens significan más gradiente, así que
mejora más rápido que los demás; al mejorar más rápido, gana aún más tokens en
el siguiente lote. Lleve el bucle hacia adelante y una capa construida con $E$
expertos rutea a través de dos o tres, y el resto nunca se entrena. Switch
Transformers observó capas enteras colapsadas a un único experto
\citep{fedus2021switch}. Esta es la trayectoria por defecto de un MoE sin
balancear, no un accidente raro, y es la razón de que exista la
Sección~\ref{app:moe-balancing-es}: la pérdida auxiliar y el sesgo de experto
son dos formas independientes de romper el bucle, la elección por experto
(Sección~\ref{app:moe-routing-es}) lo elimina por construcción y las capas
hash lo eliminan fijando la selección de raíz.

\paragraph{El rezagado.}
Incluso una capa que balancea bien en promedio fluctúa por paso: un lote puede
enviar un tercio de sus tokens a un solo experto. Siguen dos costos. El paso no
puede terminar hasta que termina el experto más ocupado, así que un único
experto caliente fija el tiempo de todo, y en entrenamiento
distribuido---donde distintos expertos viven en distintos dispositivos---el
desbalance se vuelve filas de dispositivos ociosos esperando a uno. Y un experto
que recibe un ráfaga necesita memoria de búfer proporcional a la ráfaga. El
factor de capacidad (Sección~\ref{app:moe-capacity-es}) acota ambos limitando
cuántas asignaciones debe honrar un experto. El indicador de ambos modos de
falla es MaxVio (Sección~\ref{app:moe-balancing-es}): el colapso y el rezagado
se leen ambos como el experto más ocupado despegándose de la media.

\subsection{Puntuación del router: de logits a pesos}
\label{app:moe-router-es}

El router es una única capa lineal, $z_t = W_r x_t \in \mathbb{R}^{E}$ con
$W_r \in \mathbb{R}^{E \times d}$: un logit por experto. Una función de
compuerta convierte los $E$ logits en $E$ puntuaciones, y el top-$k$ de las
puntuaciones selecciona los expertos. Conviene separar los dos oficios de la
función de compuerta: \emph{ordena} (qué expertos se seleccionan) y
\emph{pondera} (cuánto contribuye cada experto seleccionado). La implementación
soporta las tres funciones de compuerta usadas por los modelos abiertos
actuales:
\begin{align}
  \text{softmax}:&\quad s_{t,e} = \frac{e^{z_{t,e}}}{\sum_{e'} e^{z_{t,e'}}}, \\
  \text{sigmoid}:&\quad s_{t,e} = \sigma(z_{t,e}), \\
  \text{sqrt-softplus}:&\quad s_{t,e} = \sqrt{\operatorname{softplus}(z_{t,e})}.
\end{align}
Las tres son monótonas en $z_{t,e}$---un logit mayor siempre da una puntuación
mayor---y las funciones monótonas nunca reordenan. Así que para un $z_t$ dado
las tres seleccionan \emph{exactamente los mismos expertos}; la elección de la
función de compuerta es una elección sobre los pesos, no sobre la selección.
En palabras:
\begin{itemize}[leftmargin=1.4em]
  \item \textbf{softmax} hace competir a los expertos por un pastel fijo. Las
        puntuaciones suman uno, así que la puntuación del experto $e$ depende
        de todos los demás logits a través del denominador: subir un logit
        baja a todos los otros. A medida que los logits crecen, la
        distribución se afila hacia el ganador se lleva todo.
  \item \textbf{sigmoid} juzga a cada experto aislado. $\sigma(z_{t,e})$
        depende solo de su propio logit, así que no hay competencia: todos los
        expertos pueden puntuar bajo a la vez (los pesos seleccionados suman
        menos de uno y la rama ruteada se encoge) o alto a la vez. La
        renormalización y el factor de escala de abajo existen en gran parte
        para gestionar justamente esa libertad \citep{deepseek2024v3}.
  \item \textbf{sqrt-softplus} queda entre las dos: como sigmoid es por
        experto, pero $\sqrt{\operatorname{softplus}(z)}$ decae más suave por
        el lado negativo y nunca se satura en exactamente cero
        \citep{deepseek2026v4}.
\end{itemize}
\texttt{router.scoring} elige la función (por defecto \texttt{softmax}).

\paragraph{Escala de salida.}
Dos controles fijan la magnitud de la rama ruteada, y solo tienen sentido en
pareja. Con \texttt{router.renormalize} (por defecto verdadero), las
puntuaciones seleccionadas se reescalan para sumar uno,
\begin{equation}
  \hat w_{t,j} = \frac{s_{t,i_{t,j}}}{\sum_{l=1}^{k} s_{t,i_{t,l}}},
\end{equation}
de modo que la rama ruteada pasa a ser un \emph{promedio} ponderado de $k$
salidas de expertos y no una suma. Promediar encoge. Tome $k = 4$ expertos
cuyas salidas tienen cada una valor eficaz $1$ y son incorrelacionadas, con
pesos casi iguales de $\approx 1/4$: la varianza del promedio es
$\tfrac{1}{4^2}(1+1+1+1) = \tfrac{1}{4}$, así que su RMS es
$1/\sqrt{k} = 0.5$. El promedio mide la mitad que la respuesta de cualquier
experto individual, y el encogimiento empeora a medida que crece $k$.
Renormalizar aun así se paga solo---elimina la variación entre tokens de la
masa de probabilidad top-$k$ y redujo la pérdida de validación en el estudio
de granularidad de \citet{krajewski2025scaling}---así que el factor de escala
ruteado $g$ (\texttt{router.routed\_scaling\_factor}, por defecto $1$) existe
para deshacer el encogimiento: $g \approx \sqrt{k}$ recupera el RMS de un
experto individual. Los valores publicados coinciden con la derivación. Kimi
K2 usa $2.827 \approx \sqrt{8}$; DeepSeek-V3 usa $g = 2.5$ con $k = 8$,
levemente por debajo de $\sqrt{8} \approx 2.83$ porque las salidas reales de
los expertos están parcialmente correlacionadas y se cancelan menos que unas
independientes \citep{deepseek2024v3}; y Kimi K3 usa $g = 1.0$ con $k = 16$ y
dos expertos compartidos que cargan parte de la rama \citep{kimiteam2026k3}.

\paragraph{Pre-softmax.}
Con \texttt{router.pre\_softmax} el softmax se calcula sobre los $E$ expertos
\emph{antes} del corte top-$k$, de modo que los pesos conservados son
probabilidades crudas de la distribución completa y no se renormalizan. Para
ver por qué el ruteo softmax con $k{=}1$ lo necesita, recorra el camino por
defecto paso a paso: seleccionar primero y después aplicar softmax solo sobre
los logits \emph{seleccionados}. Con un único logit seleccionado ese softmax
es $e^{z}/e^{z} = 1$---siempre, sin importar qué tan bien matchee el experto
con el token. Un peso que vale constantemente uno tiene derivada cero, así
que el router no recibiría gradiente y nunca podría aprender a qué experto
preferir: la selección quedaría congelada en su inicialización. Pre-softmax
lo repara: el peso conservado es la probabilidad máxima del softmax completo,
un número genuino en $(0,1)$ que se mueve cuando se mueve cualquier logit, y
por ahí fluye el gradiente. La capa de configuración detecta esa combinación,
activa \texttt{pre\_softmax} automáticamente y emite una advertencia en vez
de rechazar configuraciones heredadas.

\paragraph{Agudeza y precisión.}
Tres mecanismos controlan cuán puntiagudo, y numéricamente confiable, es el
router. El \emph{jitter de entrada} (\texttt{router.jitter\_eps}) multiplica
el token por $x \cdot (1 + \epsilon(2u-1))$ con $u$ uniforme antes de la
proyección: ruido multiplicativo aleatorio que penaliza decisiones de ruteo
construidas sobre direcciones frágiles de $x_t$, un regularizador con el que
Switch fue entrenada \citep{fedus2021switch}. La \emph{precisión del router}
(\texttt{router.dtype}) es \texttt{float32} por defecto; el router es diminuto
($d \times E$ parámetros), así que la precisión completa no cuesta nada, y lo
que compra es estabilidad---un router bf16 divergió en Switch Transformers y
en baja precisión los empates del top-$k$ se rompen arbitrariamente
\citep{fedus2021switch}. Finalmente, la \emph{z-loss} de ST-MoE penaliza la
deriva de logits. El log-sum-exp es un máximo suave de los logits (grande
siempre que \emph{algún} logit sea grande), así que
\begin{equation}
  \mathcal{L}_z = \beta\,\frac{1}{T}\sum_{t=1}^{T}
  \Bigl(\log \sum_{e=1}^{E} e^{z_{t,e}}\Bigr)^{2},
  \label{eq:zloss-es}
\end{equation}
crece cuadráticamente siempre que cualquier logit del router crezca en
magnitud. Minimizarla mantiene los logits pequeños, el softmax sin saturar y
al router lejos de la sobreconfianza \citep{zoph2022stmoe}. Aquí $\beta$ es
\texttt{router.z\_loss\_coeff} (por defecto $0$; ST-MoE usa $10^{-3}$); la
pérdida se devuelve como \texttt{moe\_router\_z\_loss} y se suma a la pérdida
de tarea.

\subsection{Estrategia de ruteo: quién selecciona y cómo}
\label{app:moe-routing-es}

La Sección~\ref{app:moe-router-es} fijó \emph{cuánto} pesa cada experto
elegido; esta fija \emph{quién elige a quién}. La literatura tiene tres
respuestas, y cada una tiene una consecuencia estructural.

\paragraph{Elección por token.}
La opción por defecto (\texttt{routing.type: token\_choice}): cada token
conserva sus propios $k$ mejores expertos bajo $s_t + b$, donde $b$ es el sesgo
de balanceo de la Sección~\ref{app:moe-balancing-es}. La consecuencia es que
nadie garantiza trabajo a ningún experto: un experto puede ganar todos los
tokens o ninguno, que es exactamente el bucle de colapso de la
Sección~\ref{app:moe-failures-es}. La elección por token sale por eso
acompañada de un mecanismo de balanceo. Es la estrategia entrenada por
Megatron y usada por la mayoría de los modelos abiertos
\citep{fedus2021switch,deepseek2024v3}.

\paragraph{Elección por experto.}
Invierta la dirección de la decisión: con
\texttt{routing.type: expert\_choice}, cada experto conserva sus mejores
$Tk/E$ tokens \citep{zhou2022expertchoice}. Cada experto recibe exactamente
el mismo número de tokens, así que las cargas son iguales por construcción y
la maquinaria de balanceo no es solo innecesaria sino contradictoria---el
esquema prohíbe combinar elección por experto con una pérdida auxiliar o un
sesgo de experto. El precio es la causalidad: si a un token se le sirve
depende de los \emph{otros} tokens del lote, incluidos los posteriores. Eso
es inofensivo en un encoder, donde toda la secuencia está disponible a la
vez, y fatal en decodificación autorregresiva, donde el token $t$ debe
rutearse antes de que existan los tokens $t{+}1, t{+}2, \dots$; el ruteo de
entrenamiento y el de inferencia discreparían por construcción. De ahí la
recomendación: experimentos de encoder y de entrenamiento sí, serving
autorregresivo no.

\paragraph{Capas hash.}
Prescinda de la decisión aprendida por completo: las primeras
\texttt{routing.hash\_layers} capas MoE seleccionan expertos desde una tabla
fija $\mathrm{tid} \rightarrow (e_1,\dots,e_k)$ \citep{roller2021hash}, como
hace DeepSeek-V4 en sus tres primeras capas MoE \citep{deepseek2026v4}. La
observación detrás: en las primeras capas los tokens siguen cerca de sus
embeddings y la decisión de ruteo es en su mayoría identidad de
superficie---qué token es, no qué significa---y una tabla de búsqueda captura
eso. La tabla nunca cambia, así que no puede colapsar. El router aprendido
sigue produciendo los pesos de mezcla, así que sigue entrenándose a través de
ellos; simplemente no recibe gradiente de la selección misma. La tabla se
construye con una asignación aleatoria con semilla
(\texttt{hash\_table: random}) o un ciclo round-robin (\texttt{round\_robin});
Roller et al.\ hallaron que las tablas aleatorias y las voraces por frecuencia
alcanzan perplejidad similar, con mejor balance la voraz \citep{roller2021hash}.
Como una tabla fija nunca puede rebalancearse, el esquema prohíbe combinarla
con \texttt{balancing.bias\_update}.

\subsection{Balanceo de carga}
\label{app:moe-balancing-es}

La Sección~\ref{app:moe-failures-es} describió el bucle de colapso: un experto
levemente adelante recibe más tokens, entrena más rápido y se adelanta más. Se
implementan dos remedios independientes; cualquiera funciona solo, y
DeepSeek-V3 combina ambos \citep{deepseek2024v3}. Difieren en \emph{cómo}
alcanzan el ruteo: la pérdida auxiliar trabaja a través del gradiente, el sesgo
trabaja por fuera de él.

\paragraph{Pérdida auxiliar.}
El obstáculo para arreglar el desbalance directamente es que lo que hay que
arreglar---cuántas veces fue elegido cada experto---es un conteo, y los conteos
no tienen derivada: ser elegido $500$ veces en vez de $499$ no es una operación
que la retropropagación pueda influir. El sustituto de Switch
\citep{fedus2021switch} multiplica la fracción de asignaciones de cada experto
por su probabilidad media de router,
\begin{equation}
  \mathcal{L}_{\text{aux}} = \alpha\,E \sum_{e=1}^{E} f_e P_e,
  \qquad
  f_e = \frac{1}{Tk}\sum_{t,j}\mathbb{1}[i_{t,j}=e],
  \qquad
  P_e = \frac{1}{T}\sum_{t} s_{t,e},
  \label{eq:aux-es}
\end{equation}
donde $f_e$ (la fracción de las $Tk$ asignaciones) carga el desbalance pero no
el gradiente, y $P_e$ (la puntuación media) sí carga gradiente. El producto es
el truco: bajar la pérdida significa bajar $P_e$ precisamente donde $f_e$ es
alto, así que el optimizador presiona las probabilidades del router exactamente
por los expertos sobrecargados, y las selecciones futuras se reparten. Con
puntuación softmax las puntuaciones suman uno por token y el mínimo está en la
distribución uniforme, $f_e = k/E$ y $P_e = 1/E$; el prefactor $E$ hace entonces
que la pérdida en el mínimo no dependa del número de expertos (exactamente
$\alpha$ con $k{=}1$), así que un mismo coeficiente significa lo mismo con
$E{=}8$ y con $E{=}256$. El coeficiente mismo es un compromiso: demasiado
grande pelea contra la pérdida de lenguaje, demasiado pequeño deja a MaxVio
alto \citep{wang2024auxfree}; el valor por defecto es $0.01$, cercano a Switch
y DeepSeekMoE \citep{fedus2021switch,dai2024deepseekmoe}. El alcance se elige
con \texttt{balancing.aux\_loss}: \texttt{batch} (un micro-batch, por
defecto), \texttt{sequence} (una pérdida por secuencia, la guarda de
DeepSeek-V3 contra que una sola secuencia se apile en pocos expertos---balancee
el batch y aun así un documento largo puede monopolizar sus expertos
favoritos) o \texttt{none}. Balancear por secuencia o por micro-batch puede
impedir que se formen expertos de dominio, mientras que balancear sobre el
batch global deja desigual la carga por dispositivo \citep{qiu2025demons}; el
runtime de un solo dispositivo implementa los dos primeros alcances y
documenta el global como trabajo futuro.

\paragraph{Sesgo libre de pérdida auxiliar.}
El segundo remedio ignora los gradientes por completo y funciona como un
termostato: mida la temperatura, mueva la perilla, repita. Con
\texttt{balancing.bias\_update} en \texttt{sign} o \texttt{centered\_sign},
cada experto lleva un sesgo $b_e$ que se suma a las puntuaciones \emph{solo
antes de la selección---nunca a los pesos de mezcla}---y se mueve una vez por
paso de optimizador:
\begin{equation}
  b_e \leftarrow b_e + \gamma\,\operatorname{sign}\!\left(\bar c - c_e\right),
  \qquad
  \bar c = \frac{1}{E}\sum_{e'} c_{e'},
  \label{eq:bias-es}
\end{equation}
donde $c_e$ cuenta las asignaciones del paso y $\gamma$ es
\texttt{balancing.bias\_update\_rate} (por defecto $10^{-3}$). A un experto por
encima de la carga media ($c_e > \bar c$) se le baja el sesgo en $\gamma$, así
que gana menos tokens en el siguiente paso; a uno por debajo se le sube. Como
el sesgo nunca toca los pesos de mezcla, redirige tráfico sin cambiar lo que
aportan los expertos seleccionados---por eso no cuesta pérdida de lenguaje, a
diferencia de la pérdida auxiliar \citep{wang2024auxfree}.

Por qué el esquema restringe el sesgo a la puntuación sigmoid o sqrt-softplus:
bajo softmax cada puntuación depende de cada logit a través del denominador,
así que mover el logit de un experto reescala las $E$ puntuaciones a la vez y
el sesgo dirigiría la selección \emph{y} distorsionaría todos los pesos de
mezcla simultáneamente. Bajo sigmoid o sqrt-softplus cada puntuación es
función de su propio logit nada más, así que $b_e$ mueve al experto $e$ en el
ranking y no toca nada más---el direccionamiento queda puro.

La variante \texttt{centered\_sign} resta el paso medio
($\tfrac{1}{E}\sum_{e}\operatorname{sign}(\bar c - c_e) = 0$ en aritmética
exacta, no necesariamente en precisión reducida) para que la media del sesgo no
derive, lo que importa cuando el sesgo se guarda en precisión reducida
\citep{wang2024auxfree}. El sesgo es un buffer fp32 no persistente: es estado
de entrenamiento, nunca un parámetro de inferencia, y \texttt{eval()} lo
congela. La actualización la aplica \texttt{update\_moe\_expert\_bias()}
después de \texttt{optimizer.step()} en el entrenador.

\paragraph{MaxVio.}
El balance se lee con la violación máxima,
\begin{equation}
  \mathrm{MaxVio} = \frac{\max_e c_e - \bar c}{\bar c},
\end{equation}
un único número por capa, registrado como \texttt{moe/max\_vio}. Cómo leerlo:
$0$ es perfectamente parejo; $0.2$--$0.5$ es fluctuación ordinaria; $1$
significa que el experto más ocupado recibe el doble de la carga media
\citep{wang2024auxfree}. El valor absoluto importa menos que la tendencia---un
MaxVio que sube a lo largo de los pasos es el bucle de colapso de la
Sección~\ref{app:moe-failures-es} corriendo en tiempo real.

\subsection{Expertos: cantidad, ancho, espacio latente y compartición}
\label{app:moe-experts-es}

\paragraph{Cantidad y ancho.}
Los dos presupuestos de la Sección~\ref{app:moe-overview-es} se gastan aquí por
separado: \texttt{num\_experts} $\times$ \texttt{expert\_hidden\_size} es
(aproximadamente) el total de parámetros ruteados---lo que paga la
memoria---mientras que \texttt{top\_k} $\times$ \texttt{expert\_hidden\_size}
es el cómputo por token---lo que paga el usuario.
\texttt{expert\_hidden\_size} fija el ancho intermedio de cada experto y cae a
\texttt{ffn.hidden\_size} cuando no se define. Mantener los dos anchos
separados es esencial para las arquitecturas reales: DeepSeek-V3 tiene un ancho
denso de $18432$ pero expertos de ancho $2048$ \citep{deepseek2024v3}---sus
capas densas son nueve veces más anchas que sus expertos. La granularidad es el
refinamiento del mismo compromiso: con un presupuesto fijo de parámetros,
``$16$ expertos de ancho $4096$'' y ``$64$ expertos de ancho $1024$''
almacenan lo mismo. La división más fina deja a cada experto especializarse más
estrechamente, lo que mejoró la pérdida hasta $128$ expertos de ancho $d/4$ con
top-8 en un estudio \citep{liew2026principled}; pero los expertos más angostos
también ven menos tokens cada uno por paso, y al contabilizar el throughput de
entrenamiento se predice que los mismos total y activos favorecen una
configuración \emph{menos} dispersa \citep{krajewski2025scaling}. El resumen
honesto: más fino ayuda con FLOPs fijos, y la evidencia con reloj fijo es
mixta.

\paragraph{MoE latente.}
Envíe una copia comprimida en lugar del token completo. LatentMoE envuelve a
los expertos entre una proyección de bajada y una de subida: cada token se
proyecta de $d$ a un ancho latente $\ell$ antes del despacho, y el resultado
combinado se proyecta de vuelta después \citep{elango2026latentmoe},
\begin{equation}
  \ell_t = W_\downarrow x_t,\qquad
  y_t = W_\uparrow \sum_{j=1}^{k} \hat w_{t,j} E_{i_{t,j}}(\ell_t).
  \label{eq:latent-es}
\end{equation}
Cada una de las $k$ copias despachadas lleva ahora $\ell$ números en vez de
$d$. Con razón de compresión $\alpha = d/\ell$, mover y procesar una copia
cuesta $\alpha$ veces menos, así que los mismos parámetros, cómputo y tráfico
de expertos permiten $\alpha$ veces más expertos y un top-$k$ mayor. La apuesta
es que los expertos pueden trabajar sobre un boceto con pérdida del token;
NVIDIA reporta que aguanta hasta $\alpha \leqslant 4$ ($+5.65$ puntos de
MMLU-Pro a 95B totales con $\alpha = 4$) \citep{elango2026latentmoe}. Kimi K3
agrega un RMSNorm antes de la proyección de subida para detener el crecimiento
de activaciones de la rama ruteada a 2.8T parámetros, expuesto como
\texttt{latent.norm\_before\_up\_proj} \citep{kimiteam2026k3}.

\paragraph{Expertos compartidos.}
Una capa ruteada no garantiza que el conocimiento común quede representado:
cada experto aprende lo que sus visitantes necesitan, así que $E$ expertos
pueden re-aprender cada uno la misma sintaxis y puntuación. Los expertos
siempre activos de \texttt{shared.count} corren fuera del router sobre
\emph{cada} token---los médicos generalistas de la clínica de la
Sección~\ref{app:moe-overview-es}---absorbiendo el conocimiento que todo token
necesita y liberando a los especialistas ruteados para gastar su capacidad en
lo que difiere entre tokens \citep{dai2024deepseekmoe}. Su ancho es
\texttt{shared.hidden\_size} (por defecto, el ancho ruteado). Se suman con peso
uno, o multiplicados por una compuerta aprendida por token
$\sigma(w_s^{\top} x_t) \in (0,1)$ cuando \texttt{shared.gate: sigmoid}
(Qwen3-Next) \citep{qwen2025qwen3}; la forma de peso fijo coincide con
DeepSeek-V3 \citep{deepseek2024v3}. No son gratis: un experto compartido corre
sobre cada token, así que el cómputo por token crece con su ancho, y las
ganancias de calidad medidas son pequeñas e inconsistentes
\citep{li2026slicing}. Agregue uno cuando de otro modo subiría
\texttt{top\_k}, no encima de él.

\paragraph{Tope de activación.}
Una patología rara pero real: expertos individuales cuyas normas de salida
crecen mientras la pérdida queda plana---Step 3.5 Flash encontró esos atípicos
y el recorte los contuvo sin cambiar la pérdida \citep{huang2026step35}.
\texttt{experts.activation\_clamp} acota el intermedio con compuerta,
\begin{equation}
  \mathrm{GatedFFN}(x) = W_{\text{down}}\!\left(
  \phi\bigl(\operatorname{clamp}(xW_{\text{gate}}, -c, c)\bigr) \odot (xW_{\text{up}})
  \right),
\end{equation}
con $c$ = \texttt{activation\_clamp} ($0$ lo desactiva). DeepSeek-V4.1 usa
$c = 10$ y Step 5 Preview $c = 7$ \citep{deepseek2026v4flash}. Trátelo como un
seguro contra expertos atípicos, no como un valor por defecto: existe para que
la señal de monitoreo de la Sección~\ref{app:moe-monitoring-es} (razón de
normas de expertos creciendo) tenga un remedio barato. El tope suave exacto de
SiTU-GLU que usa Kimi K3 se aproxima aquí con un recorte de la entrada de la
compuerta y se documenta como desviación \citep{kimiteam2026k3}.

\subsection{Disposición por capas}
\label{app:moe-layout-es}

No toda capa necesita ser MoE. Las primeras capas de un transformer computan
características genéricas que todo token necesita sin importar su contenido; la
especialización paga después. \texttt{layout.dense\_layers} mantiene densas las
primeras $N$ capas (DeepSeek-V3 y GLM conservan tres), y
\texttt{layout.every\_n\_layers} intercala capas densas y MoE (Llama 4
Maverick hace MoE una de cada dos capas)
\citep{deepseek2024v3,meta2025llama4}. Las capas densas usan
\texttt{ffn.hidden\_size}; las MoE usan \texttt{expert\_hidden\_size}. La
configuración del modelo expone \texttt{is\_moe\_layer(i)},
\texttt{moe\_layer\_mask()} y \texttt{moe\_layer\_count()}; las pilas de
encoder, decoder, ViT y audio construyen una FFN ruteada o densa por capa
física. Las capas densas iniciales reducen el número de capas MoE a bajo costo
con FLOPs igualados \citep{krajewski2025scaling}, pero no corrigen por sí
solas el desbalance de las primeras capas: las primeras capas MoE tras el
prefijo siguen siendo las más propensas al colapso y las que hay que
monitorear.

\subsection{Capacidad y descarte de tokens}
\label{app:moe-capacity-es}

Por defecto el ruteo es dropless: un experto acepta todos los tokens que el
router le envía (\texttt{routing.capacity\_factor: 0}) y las cargas por paso
fluctúan libremente. Un factor positivo $\kappa$ limita cada experto a
\begin{equation}
  C = \left\lceil \kappa\,\frac{Tk}{E} \right\rceil
\end{equation}
asignaciones. La fórmula es la carga media con holgura: $T$ tokens hacen $k$
solicitudes cada uno, dando $Tk$ asignaciones repartidas entre $E$ expertos,
así que la media es $Tk/E$ y $\kappa$ es el margen. Ejemplo trabajado:
$T = 4096$ tokens, $E = 8$, $k = 2$ da $8192$ asignaciones y una media de
$1024$ por experto; $\kappa = 1.5$ fija $C = 1536$. Un experto que recibe
$2000$ solicitudes atiende sus $1536$ de mayor puntuación y descarta las
otras $464$: esos tokens salen de la capa sin la contribución de ese experto
(su peso se pone en cero), y \texttt{dropped\_fraction} registra la fracción
de asignaciones eliminadas ($464$ de $8192 \approx 5.7\%$ en el ejemplo). Lo
que compra el tope es el remedio al rezagado de la
Sección~\ref{app:moe-failures-es}: un límite duro al búfer de cada experto
(memoria) y, en entrenamiento distribuido, la garantía de que ningún experto
caliente fija el tiempo del paso. El costo es pérdida de calidad silenciosa:
Pangu Ultra MoE descartó cerca de 6--8\% de tokens con $\kappa = 1.5$ y su
corrida dropless mantuvo menor pérdida, con una brecha que crecía con el
número de expertos \citep{pangu2025ultra}, y MegaBlocks alcanzó su menor
pérdida cuando no se descartó ningún token \citep{gale2022megablocks}. Dropless
sigue siendo el valor por defecto y la opción recomendada para fine-tuning,
donde los tokens descartados pueden caer dentro de respuestas de instrucción.

\subsection{Monitoreo}
\label{app:moe-monitoring-es}

Cada \texttt{MoELayer} expone \texttt{last\_auxiliary\_losses} (las pérdidas
\texttt{moe\_load\_balancing\_loss} y \texttt{moe\_router\_z\_loss} ponderadas
por sus coeficientes) y \texttt{last\_moe\_stats}. El modelo las agrega sobre
las capas mediante \texttt{aggregate\_moe\_aux\_losses()} y suma las pérdidas a
la de tarea; el entrenador registra los escalares en el callback de métricas y
en Weights \& Biases. Léase como cinco preguntas, una señal cada una:
\begin{itemize}[leftmargin=1.4em]
  \item \textbf{¿Está pareja la carga?} MaxVio (\texttt{moe/max\_vio}) por
        capa: cerca de $0$ sano, $\geqslant 1$ severo
        (Sección~\ref{app:moe-balancing-es}); un valor que sube a lo largo de
        los pasos es el bucle de colapso en marcha.
  \item \textbf{¿Se comprometió el router?} Entropía de ruteo y peso top-1.
        Una compuerta plana que apenas promedia sus $k$ expertos se ve como
        entropía cercana a $\log k$ y peso top-1 cercano a $1/k$ (sin
        especialización); una compuerta de ganador se lleva todo muestra
        entropía cercana a $0$ y top-1 cercano a $1$. Un entrenamiento sano
        deriva del primer estado hacia un ruteo comprometido y estructurado
        sin llegar al segundo extremo.
  \item \textbf{¿Están sanos los expertos?} La razón de normas de expertos:
        la mayor norma de salida por experto sobre la mediana de la capa. Muy
        por encima de uno, un experto atípico dominando la rama---la patología
        que contiene el tope de activación; clavada en cero, un experto muerto
        que nadie selecciona \citep{huang2026step35}.
  \item \textbf{¿Con cuánta fuerza tira el balanceo y está derivando el
        router?} La pérdida de balanceo y la z-loss (coeficientes ya divididos
        en las estadísticas registradas): la primera debe relajarse a medida
        que mejora el balance, la segunda sigue la deriva de logits.
  \item \textbf{¿Está apretando la capacidad?} La fracción descartada:
        proporción de asignaciones eliminadas por el factor de capacidad; es
        $0$ salvo que \texttt{routing.capacity\_factor} sea positivo.
\end{itemize}

Cinco reglas prácticas cierran la rampa, cada una consecuencia de las secciones
anteriores y no maquinaria nueva:
\begin{enumerate}[leftmargin=1.4em]
  \item Arranque con los valores por defecto: ruteo dropless, pérdida auxiliar
        de alcance batch en $0.01$, sin sesgo. Son la configuración a la que
        convergen la mayoría de los presets tras sus propias desviaciones.
  \item Si el problema es el balance, vaya primero al sesgo antes que a subir
        el coeficiente auxiliar: el sesgo redirige tráfico solo a través de la
        selección y no cuesta pérdida de lenguaje \citep{wang2024auxfree},
        mientras que un $\alpha$ mayor pelea contra la pérdida de tarea.
  \item Vigile MaxVio en las primeras capas MoE tras cualquier prefijo denso;
        ahí empieza el colapso (Sección~\ref{app:moe-layout-es}).
  \item Softmax con $k{=}1$ necesita \texttt{pre\_softmax} o el router nunca
        aprende (Sección~\ref{app:moe-router-es}); el cargador lo repara, pero
        la advertencia vale la pena entenderla.
  \item Presupueste el cómputo como
        $\texttt{top\_k} \times \texttt{expert\_hidden\_size} +
        \texttt{shared}$ antes de habilitar: los expertos compartidos y el
        ancho latente cambian la factura por token, y el remedio más barato
        para expertos lentos es un presupuesto menor, no más maquinaria.
\end{enumerate}

\subsection{Contrato de esquema}
\label{app:moe-schema-es}

La Tabla~\ref{tab:moe-schema-es} lista cada clave \texttt{model.moe.*}, su
valor por defecto, qué controla y la referencia que la gobierna. El esquema
JSON es el contrato ejecutable; la validación en runtime agrega las reglas
entre claves: $top\_k \leqslant num\_experts$, softmax con $k=1$ implica
pre-softmax, el sesgo requiere puntuación sigmoid o sqrt-softplus, la elección
por experto es incompatible con los términos de balanceo, las capas hash
requieren elección por token y prohíben el sesgo, el ancho latente debe estar
en $(0, d)$ y el número de capas debe ser positivo.

{\small
\begin{longtable}{@{}L{0.27\linewidth}L{0.13\linewidth}L{0.32\linewidth}L{0.20\linewidth}@{}}
\caption{Contrato de esquema MoE (\texttt{model.moe.*}).}\label{tab:moe-schema-es}\\
\toprule
\textbf{Clave} & \textbf{Defecto} & \textbf{Controla} & \textbf{Referencia} \\
\midrule
\endfirsthead
\toprule
\textbf{Clave} & \textbf{Defecto} & \textbf{Controla} & \textbf{Referencia} \\
\midrule
\endhead
\texttt{enabled} & \texttt{true} & FFN ruteada vs.\ densa & \citet{shazeer2017moe} \\
\texttt{num\_experts} & 8 & tamaño del pool ruteado & \citet{fedus2021switch} \\
\texttt{top\_k} & 2 & expertos usados por token & \citet{fedus2021switch} \\
\texttt{expert\_hidden\_size} & \texttt{null} & ancho del experto (cae a \texttt{ffn.hidden\_size}) & \citet{deepseek2024v3} \\
\texttt{router.scoring} & \texttt{softmax} & función logits$\to$puntuaciones & \citet{deepseek2024v3,deepseek2026v4} \\
\texttt{router.pre\_softmax} & \texttt{false} & softmax antes del corte top-$k$ & \citet{fedus2021switch} \\
\texttt{router.renormalize} & \texttt{true} & pesos seleccionados suman uno & \citet{krajewski2025scaling} \\
\texttt{router.routed\_scaling\_factor} & 1.0 & magnitud de la rama ruteada $g$ & \citet{deepseek2024v3,kimiteam2026k3} \\
\texttt{router.dtype} & \texttt{float32} & precisión de cómputo del router & \citet{fedus2021switch,zoph2022stmoe} \\
\texttt{router.jitter\_eps} & 0.0 & fuerza del jitter de entrada & \citet{fedus2021switch} \\
\texttt{router.z\_loss\_coeff} & 0.0 & penalización de deriva de logits $\beta$ & \citet{zoph2022stmoe} \\
\texttt{routing.type} & \texttt{token\_choice} & quién elige a quién & \citet{zhou2022expertchoice} \\
\texttt{routing.capacity\_factor} & 0.0 & tope por experto $\kappa$ (0: dropless) & \citet{pangu2025ultra} \\
\texttt{routing.hash\_layers} & 0 & primeras $N$ capas ruteadas por tabla & \citet{roller2021hash,deepseek2026v4} \\
\texttt{routing.hash\_table} & \texttt{random} & construcción de la tabla & \citet{roller2021hash} \\
\texttt{routing.hash\_seed} & 42 & semilla de la tabla & \citet{roller2021hash} \\
\texttt{balancing.aux\_loss} & \texttt{batch} & alcance de la pérdida & \citet{dai2024deepseekmoe,deepseek2024v3} \\
\texttt{balancing.aux\_loss\_coeff} & 0.01 & peso de la pérdida $\alpha$ & \citet{deepseek2024v3} \\
\texttt{balancing.bias\_update} & \texttt{none} & termostato off / \texttt{sign} / centrado & \citet{wang2024auxfree} \\
\texttt{balancing.bias\_update\_rate} & 0.001 & paso del termostato $\gamma$ & \citet{wang2024auxfree} \\
\texttt{shared.count} & 0 & expertos siempre activos & \citet{dai2024deepseekmoe} \\
\texttt{shared.hidden\_size} & \texttt{null} & ancho del experto compartido & \citet{dai2024deepseekmoe} \\
\texttt{shared.gate} & \texttt{none} & peso fijo o compuerta aprendida & \citet{qwen2025qwen3} \\
\texttt{latent.enabled} & \texttt{false} & comprimir antes del despacho & \citet{elango2026latentmoe} \\
\texttt{latent.dim} & \texttt{null} & ancho latente $\ell$ & \citet{elango2026latentmoe} \\
\texttt{latent.norm\_before\_up\_proj} & \texttt{false} & RMSNorm antes de $W_\uparrow$ & \citet{kimiteam2026k3} \\
\texttt{layout.dense\_layers} & 0 & prefijo denso $N$ & \citet{deepseek2024v3} \\
\texttt{layout.every\_n\_layers} & 1 & MoE cada $n$-ésima capa & \citet{meta2025llama4} \\
\texttt{experts.activation\_clamp} & 0.0 & recorte de la entrada de compuerta $c$ (0: off) & \citet{huang2026step35} \\
\bottomrule
\end{longtable}
}

\subsection{Presets con nombre}
\label{app:moe-presets-es}

Los presets de \texttt{configs/} reflejan las configuraciones publicadas, con
las desviaciones documentadas para un solo dispositivo. La
Tabla~\ref{tab:moe-presets-es} resume el bloque MoE de cada preset MoE; los
valores se verificaron contra el \texttt{config.json} de Hugging Face de cada
modelo cuando estaba disponible. Leída contra la
Sección~\ref{app:moe-failures-es}: todo preset frontera reciente responde al
colapso con el sesgo, la pérdida auxiliar o ambos.

{\small
\begin{longtable}{@{}L{0.18\linewidth}L{0.12\linewidth}L{0.10\linewidth}L{0.14\linewidth}L{0.12\linewidth}L{0.22\linewidth}@{}}
\caption{Presets MoE y sus bloques MoE de arquitectura real.}\label{tab:moe-presets-es}\\
\toprule
\textbf{Preset} & \textbf{Expertos / top-k} & \textbf{Ancho} & \textbf{Router} & \textbf{Compartidos} & \textbf{Balanceo / capas} \\
\midrule
\endfirsthead
\toprule
\textbf{Preset} & \textbf{Expertos / top-k} & \textbf{Ancho} & \textbf{Router} & \textbf{Compartidos} & \textbf{Balanceo / capas} \\
\midrule
\endhead
DeepSeek-V2 & 160 / 6 & 1536 & softmax, $g{=}16$ & 2 & aux-secuencia $10^{-3}$; 1 densa \\
DeepSeek-V3 & 256 / 8 & 2048 & sigmoid, $g{=}2.5$ & 1 & sesgo + aux-secuencia $10^{-4}$; 3 densas \\
GLM-4.5/4.7 & 160 / 8 & 1536 & sigmoid, $g{=}2.5$ & 1 & sesgo; 3 densas \\
GLM-5 & 256 / 8 & 2048 & sigmoid, $g{=}2.5$ & 1 & sesgo; 3 densas \\
Kimi K2 & 384 / 8 & 2048 & sigmoid, $g{=}2.827$ & 1 & sesgo + aux-secuencia; 1 densa \\
Kimi K3 & 896 / 16 & 3072 & sigmoid, $g{=}1$ & 2 & sesgo centrado; latente 3584; 1 densa \\
gpt-oss-120b & 128 / 4 & 2880 & softmax & 0 & aux, tope 7 \\
Llama 4 Scout & 16 / 1 & 8192 & softmax, pre & 0 & todas las capas \\
Llama 4 Maverick & 128 / 1 & 8192 & softmax, pre & 0 & una de cada 2 capas \\
MiniMax-Text-01 & 32 / 2 & 9216 & sigmoid & 0 & aux $10^{-3}$ \\
MiniMax-M2 & 256 / 8 & 1536 & sigmoid & 0 & sesgo + aux $10^{-3}$ \\
Qwen3-Next-80B & 512 / 10 & 512 & softmax & 1, con compuerta & aux $10^{-3}$ \\
Qwen3.8 2.4T-A95B & 512 / 10 & 2048 & softmax & 1 & aux $10^{-3}$ \\
\bottomrule
\end{longtable}
}

\paragraph{Desviaciones de un solo dispositivo.}
El runtime es CPU-first y de un solo dispositivo, así que implementa la
semántica a nivel de modelo y deja los mecanismos distribuidos para trabajo
futuro. Las desviaciones documentadas son: la pérdida a nivel de dispositivo de
DeepSeek-V2 se aproxima con la pérdida auxiliar de batch; el balanceo por
cuantiles/PID de Kimi K3 se aproxima con la regla de signo centrado; el tope
suave de SiTU-GLU de Kimi K3 se aproxima con \texttt{activation\_clamp: 4}; el
escalado de puntuaciones en la entrada de Llama 4 no se aplica (los pesos del
router se aplican a la salida del experto); y el ruteo limitado por grupos, el
paralelismo de expertos, los despachadores de tokens, las réplicas de
expertos/EPLB y la ejecución de expertos en FP8/FP4 quedan fuera de alcance.
Las capas hash de DeepSeek-V4 se soportan con \texttt{routing.hash\_layers}.

\subsection{Interacción con el entrenamiento}
\label{app:moe-training-es}

El sesgo de experto se actualiza después del paso del optimizador a partir de
los conteos acumulados en los forwards del paso, así que la acumulación de
gradiente y el micro-batching no cambian su semántica. El router no recibe
gradiente del sesgo y el sesgo es un buffer no persistente; los checkpoints
exportados no contienen por tanto estado de entrenamiento---quedan limpios para
inferencia. Las pérdidas auxiliares son tensores ponderados por coeficiente en
\texttt{last\_auxiliary\_losses} y el entrenador las suma a la pérdida de
tarea; se excluyen de la evaluación. El router y los expertos se categorizan
como bloques \texttt{router} y \texttt{ffn} respectivamente, así que las normas
de gradiente por bloque y los conteos de parámetros coinciden con la taxonomía
del estimador. Finalmente, la capa MoE es token-wise y por tanto segura para
caché: las suites de paridad de caché KV mantienen la FFN densa porque un
top-$k$ discreto puede cambiar bajo las diferencias numéricas mínimas entre
forwards completos y con caché---una propiedad del ruteo discreto, no de la
caché.

\subsection{Sistemas relacionados y trabajo futuro}
\label{app:moe-future-es}

El diseño de este anexo sigue la anatomía del manual de MoE (router, ruteo,
balanceo, expertos, disposición, capacidad) y la literatura de sistemas que
resume. Trabajo futuro en este código: paralelismo de expertos y despachadores
all-to-all \citep{lepikhin2020gshard,fedus2021switch}; pérdidas auxiliares sobre
el batch global \citep{qiu2025demons}; balanceo por cuantiles/PID
\citep{jin2026idbalancing}; upcycling de denso a MoE y crecimiento de expertos
\citep{komatsuzaki2023upcycling}; tablas hash balanceadas por frecuencia
\citep{roller2021hash}; SiTU-GLU exacto; y bloques MoE con conexión de atajo o
paralelos que solapen el despacho con el cómputo
\citep{cai2024scmoe,longcat2025flash}.
